\begin{apendicesenv}

\chapter{CHECKLIST PRISMA 2020}
\label{apendice:prisma-checklist}

Este apêndice apresenta o checklist completo de aderência às 27 recomendações do PRISMA 2020 para relato de revisões sistemáticas. Para cada item, indica-se a localização no documento onde a informação correspondente pode ser encontrada.

\begin{longtable}{|p{0.5cm}|p{3.5cm}|p{6cm}|p{3cm}|}
\hline
\textbf{N\textordmasculine} & \textbf{Item PRISMA 2020} & \textbf{Descrição} & \textbf{Localização no PTC} \\
\hline
\endfirsthead

\multicolumn{4}{c}%
{\tablename\ \thetable\ -- Continuação da página anterior} \\
\hline
\textbf{N\textordmasculine} & \textbf{Item PRISMA 2020} & \textbf{Descrição} & \textbf{Localização no PTC} \\
\hline
\endhead

\hline \multicolumn{4}{r}{\textit{Continua na próxima página}} \\
\endfoot

\hline
\endlastfoot

%--- TÍTULO ---
1 & Título & Identificar o relatório como uma revisão sistemática & Capa, folha de rosto \\
\hline

%--- RESUMO ---
2 & Resumo estruturado & Resumo estruturado incluindo contexto, objetivo, métodos, resultados e conclusões & Resumo (pág. vi-vii) \\
\hline

%--- INTRODUÇÃO ---
3 & Justificativa & Descrever a fundamentação da revisão no contexto do conhecimento existente & Cap. 1, Seção 1.1-1.2 \\
\hline

4 & Objetivos & Fornecer declaração explícita de questões de pesquisa/objetivos & Cap. 1, Seção 1.3 \\
\hline

%--- METODOLOGIA ---
5 & Critérios de elegibilidade & Especificar critérios de inclusão/exclusão & Cap. 3, Seção 3.3 \\
\hline

6 & Fontes de informação & Especificar todas as bases consultadas e datas de cobertura & Cap. 3, Seção 3.2.1 \\
\hline

7 & Estratégia de busca & Apresentar estratégia de busca completa para pelo menos uma base & Cap. 3, Seção 3.2.2 \\
\hline

8 & Processo de seleção & Especificar processo de screening, elegibilidade e inclusão & Cap. 3, Seção 3.3 \\
\hline

9 & Avaliação de risco de viés & Especificar métodos para avaliar risco de viés nos estudos & \textit{Não aplicado -- previsto para TCC} \\
\hline

10 & Coleta de dados & Especificar métodos de extração de dados & Cap. 3, Seção 3.4 \\
\hline

11a & Dados coletados & Listar e definir todas as variáveis extraídas & Cap. 3, Seção 3.4 \\
\hline

11b & Desvios do protocolo & Relatar desvios do protocolo e justificativas & Cap. 3 (sem desvios significativos) \\
\hline

%--- RESULTADOS ---
12 & Seleção de estudos & Reportar número de estudos em cada etapa com razões para exclusões & Cap. 4, Fig. PRISMA flow \\
\hline

13 & Características dos estudos & Apresentar características relevantes dos estudos incluídos & Cap. 4, Seção 4.2 \\
\hline

14 & Risco de viés nos estudos & Apresentar avaliações de risco de viés para cada estudo & \textit{Não aplicado -- previsto para TCC} \\
\hline

15 & Resultados de estudos individuais & Apresentar resultados de cada estudo analisado & Cap. 4, Seção 4.3 \\
\hline

16a & Síntese dos resultados & Método de síntese (meta-análise ou narrativa) & Cap. 4, Seção 4.4 (narrativa) \\
\hline

16b & Métodos adicionais de síntese & Outros métodos de análise/síntese (subgrupos, sensibilidade, meta-regressão) & N/A (síntese narrativa) \\
\hline

17 & Viés de publicação & Avaliação de viés de publicação & \textit{Não aplicado -- previsto para TCC} \\
\hline

18 & Certeza da evidência & Avaliação de certeza/qualidade da evidência & \textit{Não aplicado -- previsto para TCC} \\
\hline

%--- DISCUSSÃO ---
19 & Discussão geral & Interpretação dos resultados considerando objetivos e outras evidências & Cap. 4, Discussão \\
\hline

20 & Limitações da evidência & Limitações da evidência incluída (risco de viés, inconsistência, imprecisão) & \textit{Seção explícita será adicionada no TCC} \\
\hline

21 & Limitações da revisão & Limitações dos processos da própria revisão & \textit{Seção explícita será adicionada no TCC} \\
\hline

22 & Implicações & Implicações para prática, política e pesquisa futura & Cap. 4, Conclusão \\
\hline

%--- OUTROS ---
23 & Registro e protocolo & Informações sobre registro e protocolo (PROSPERO, etc.) & Cap. 3 (justificativa de ausência) \\
\hline

24 & Suporte & Fontes de apoio financeiro e outros suportes & Folha de rosto \\
\hline

25 & Conflitos de interesse & Declarar conflitos de interesse & Folha de rosto \\
\hline

26 & Disponibilidade de dados & Declarar disponibilidade de dados, código e outros materiais & Cap. 3 (código em repositório GitHub) \\
\hline

27 & Checklist & Incluir checklist PRISMA 2020 completo & Presente neste apêndice \\
\hline

\caption{Checklist de aderência ao PRISMA 2020 \cite{PRISMA2020}}
\label{tab:prisma-checklist}
\end{longtable}

\section*{Notas Explicativas}

\textbf{Legenda:}
\begin{itemize}
    \item Itens com localização específica: totalmente atendidos no documento
    \item Itens marcados como \textit{``Não aplicado -- previsto para TCC''}: parcialmente atendidos ou previstos para fase TCC
    \item N/A: item não aplicável ao tipo/escopo desta revisão
\end{itemize}

\textbf{Itens marcados como ``previsto para TCC'':}

Os itens 9, 14, 17, 18, 20 e 21 referem-se a avaliações de qualidade metodológica, risco de viés, certeza da evidência e limitações explícitas. Estes elementos não foram aplicados na fase PTC devido a:

\begin{enumerate}
    \item \textbf{Cronograma acadêmico:} PTC focado em mapeamento exploratório
    \item \textbf{Escopo da Fase 1:} Identificação e categorização de estudos
    \item \textbf{Planejamento do TCC:} Avaliação crítica aprofundada será conduzida na Fase 2
\end{enumerate}

A aplicação do \textbf{Mixed Methods Appraisal Tool (MMAT)} \cite{MMAT2018} está programada para a fase TCC, permitindo avaliar a qualidade metodológica dos 17 estudos incluídos considerando a heterogeneidade de desenhos de pesquisa identificados.

\textbf{Referência:}

Page MJ, McKenzie JE, Bossuyt PM, et al. The PRISMA 2020 statement: an updated guideline for reporting systematic reviews. \textit{BMJ} 2021;372:n71. doi: 10.1136/bmj.n71

\end{apendicesenv}
